\rtmtight
\begin{tblr}{width=\textwidth,colspec={|Q[c,m,2.1cm]|X[l,m]|},hlines,vlines,hspan=minimal,rowsep=1.5pt,colsep=3pt,row{1}={bg=headgray,font=\bfseries}}
\SetCell{c,m}\hfil\logo\hfil & \textbf{POLITEKNIK NEGERI MALANG}\par\textbf{JURUSAN TEKNOLOGI INFORMASI}\par\textbf{PROGRAM STUDI: D4 TEKNIK INFORMATIKA} \\
\SetCell[c=2]{l} \textbf{RENCANA TUGAS MAHASISWA (RTM) DAN RUBRIK PENILAIAN} & \\
Nama Mata Kuliah & Bahasa Inggris 1 \\
Kode Mata Kuliah & RTI251005 \quad \textbf{sks / Jam perminggu:} 2 SKS/4 jam \quad \textbf{Semester:} 1 \\
Dosen Pengampu & \begin{enumerate}\item Atiqah Nurul Asri, S.Pd., M.Pd.\item Satrio Binusa Suryadi, S.S., M.Pd.\end{enumerate} \\
\end{tblr}
\assessmentblock{Tugas Menulis dan Presentasi Terstruktur}{Tugas terstruktur (tulisan dan presentasi)}{SCPMK0303-00503}{Mahasiswa mengerjakan lima tugas menulis terstruktur (instruksi instalasi software, perbandingan device, proses editing, review website, dan deskripsi dream job) masing-masing berbobot 1\% sesuai topik yang dipelajari.}{Memilih objek tugas, menulis teks dengan grammatical functions target, menyunting dengan spelling/grammar checker, menyiapkan slide/video dengan screenshot, dan mempresentasikannya.}{Dokumen teks, slide presentasi, dan/atau video dengan subtitle.}{Ketepatan grammatical functions sesuai topik & 2 & Bentuk gramatikal target (Simple Present, imperatives, sequencers, comparatives, passive, modals) digunakan benar pada hampir seluruh kalimat; kesalahan minor tidak mengganggu makna. & Bentuk gramatikal target umumnya benar, tetapi terdapat kesalahan berulang pada sebagian kalimat yang kadang mengaburkan makna. & Bentuk gramatikal target jarang digunakan atau banyak kesalahan sehingga teks sulit dipahami. \\ \\
Kelengkapan isi dan struktur teks & 1.5 & Judul, identitas, seluruh langkah/fitur, dan screenshot pendukung tersaji lengkap, runtut, dan sesuai checklist tugas. & Sebagian besar komponen tersaji, tetapi ada langkah/screenshot yang hilang atau urutan kurang runtut. & Komponen utama banyak yang hilang atau isi tidak sesuai instruksi tugas. \\ \\
Kualitas penyuntingan dan format luaran & 1.5 & Teks dan slide bebas typo, satu langkah per slide, format konsisten, dan luaran dikumpulkan tepat waktu. & Masih ada beberapa typo atau format slide kurang konsisten, tetapi luaran tetap terbaca jelas. & Banyak typo, format tidak rapi, atau luaran tidak lengkap/terlambat. \\}

\assessmentblock{Kuis Vocabulary dan Grammar (Kuis 1 dan Kuis 2)}{Kuis tertulis}{SCPMK0303-00501}{Mahasiswa mengerjakan dua kuis tertulis: Kuis 1 (Topic 1--3) pada minggu ke-6 dan Kuis 2 (Topic 4--7) pada minggu ke-14, mencakup vocabulary dan grammatical functions yang dipelajari.}{Mengerjakan soal pilihan/isian/menulis kalimat tentang vocabulary Teknik Informatika dan grammatical functions secara individu dalam waktu yang ditentukan.}{Lembar jawaban kuis tertulis (luring atau melalui LMS).}{Penguasaan vocabulary Teknik Informatika & 3 & Lebih dari 80\% kata/istilah IT dijawab atau digunakan dengan makna dan ejaan yang tepat. & 60--80\% kata/istilah IT dijawab dengan tepat; sebagian ejaan atau makna kurang akurat. & Kurang dari 60\% kata/istilah IT dijawab dengan tepat. \\ \\
Ketepatan grammatical functions & 4 & Lebih dari 80\% soal grammar (tenses, comparatives, passive, if-clause, modals) dijawab benar sesuai konteks. & 60--80\% soal grammar dijawab benar; kesalahan terjadi pada bentuk yang lebih kompleks. & Kurang dari 60\% soal grammar dijawab benar. \\ \\
Penerapan dalam kalimat/paragraf konteks TI & 3 & Kalimat/paragraf yang ditulis gramatikal, koheren, dan menggunakan istilah TI secara tepat sesuai perintah soal. & Kalimat/paragraf dapat dipahami, tetapi terdapat kesalahan struktur atau pilihan istilah yang kurang tepat. & Kalimat/paragraf tidak sesuai perintah, tidak koheren, atau penuh kesalahan struktur. \\}

\assessmentblock{Praktik Lisan: Tanya Jawab, Diskusi, Role Play, dan Presentasi}{Case Method dan praktik lisan}{SCPMK0303-00502, SCPMK0303-00503, SCPMK0303-00504}{Mahasiswa mengikuti penilaian praktik lisan berkelanjutan pada setiap topik: tanya jawab, diskusi kelompok, role play (misalnya percakapan IT Support Staff-Client), dan presentasi hasil tugas secara individu maupun kelompok.}{Memahami teks/simakan sebagai bahan, menyiapkan teks dan slide pendukung, berlatih, lalu tampil: menyapa audiens, memperkenalkan diri, menyampaikan tujuan, menjelaskan isi, menyimpulkan, dan menjawab pertanyaan.}{Penampilan lisan di kelas atau video, disertai slide/teks pendukung dan bukti latihan.}{Pemahaman listening-reading dan ketepatan respon tanya jawab & 5 & Menjawab pertanyaan tentang isi bacaan/simakan dengan benar dan menjelaskan kembali isi teks dengan kata sendiri secara akurat. & Menjawab sebagian besar pertanyaan dengan benar, tetapi penjelasan ulang isi teks kurang lengkap atau kurang akurat. & Jawaban sering keliru atau tidak mampu menjelaskan kembali isi bacaan/simakan. \\ \\
Kualitas teks dan bahan presentasi pendukung & 10 & Teks/slide lengkap sesuai checklist, bebas typo, satu ide per slide, screenshot/visual relevan, dan struktur pembuka-isi-penutup jelas. & Teks/slide mencakup isi utama, tetapi ada bagian checklist yang terlewat, typo, atau visual kurang mendukung. & Teks/slide tidak lengkap, sulit diikuti, atau tidak sesuai topik tugas. \\ \\
Kelancaran (fluency) dan pelafalan & 15 & Berbicara lancar tanpa jeda panjang, pelafalan jelas dan mudah dipahami, tidak sekadar membaca teks. & Berbicara cukup lancar dengan beberapa jeda atau pengulangan; sebagian pelafalan kurang tepat tetapi masih dipahami. & Sering berhenti lama, membaca teks kata per kata, atau pelafalan sulit dipahami. \\ \\
Ketepatan grammar dan jangkauan vocabulary lisan & 13 & Menggunakan grammatical functions target dan istilah TI secara benar dan bervariasi selama tampil. & Grammar dan vocabulary umumnya benar, tetapi variasinya terbatas atau ada kesalahan berulang. & Banyak kesalahan grammar atau vocabulary sangat terbatas sehingga pesan tidak tersampaikan. \\ \\
Organisasi presentasi dan penyelesaian tugas komunikatif & 12 & Presentasi terstruktur (greeting, tujuan, isi, kesimpulan, undangan bertanya), sesuai alokasi waktu, dan seluruh tujuan komunikatif tugas tercapai. & Struktur presentasi ada tetapi beberapa bagian terlewat, waktu kurang terkelola, atau sebagian tujuan komunikatif belum tercapai. & Presentasi tidak terstruktur, jauh dari alokasi waktu, atau tujuan komunikatif tugas tidak tercapai. \\}

\assessmentblock{UTS Bahasa Inggris 1}{Ujian Tengah Semester (tes tulis dan/atau lisan)}{SCPMK0303-00501, SCPMK0303-00502, SCPMK0303-00503}{Mahasiswa mengerjakan UTS yang mencakup vocabulary dan grammatical functions, pemahaman bacaan/simakan, serta penulisan karangan pendek berdasarkan Topic 1--3 (Computer Applications, Computer Architecture, Multimedia).}{Mengerjakan soal secara individu dalam waktu yang ditentukan: soal vocabulary-grammar, reading/listening comprehension, dan menulis karangan pendek sesuai instruksi.}{Lembar jawaban UTS (tulis) dan/atau rekaman respon lisan.}{Ketepatan vocabulary dan grammar & 5 & Lebih dari 80\% soal vocabulary dan grammar (Simple Present, imperatives, sequencers, comparatives, passive, time clauses) dijawab benar. & 60--80\% soal dijawab benar; kesalahan terkonsentrasi pada bentuk yang kompleks. & Kurang dari 60\% soal dijawab benar. \\ \\
Pemahaman bacaan/simakan & 5 & Menjawab pertanyaan comprehension dengan benar, menemukan informasi rinci, dan menyimpulkan isi teks secara tepat. & Sebagian besar pertanyaan dijawab benar, tetapi kesimpulan atau informasi rinci kurang tepat. & Jawaban comprehension banyak yang keliru atau tidak menunjukkan pemahaman isi teks. \\ \\
Kualitas karangan & 5 & Karangan sesuai perintah, terorganisasi runtut, menggunakan grammatical functions target dengan benar, dan istilah TI yang tepat. & Karangan menjawab perintah, tetapi organisasi atau ketepatan grammar-vocabulary masih lemah di beberapa bagian. & Karangan tidak sesuai perintah, tidak runtut, atau kesalahan bahasa mendominasi. \\}

\assessmentblock{UAS Bahasa Inggris 1}{Ujian Akhir Semester (tes tulis dan/atau lisan)}{SCPMK0303-00501, SCPMK0303-00502, SCPMK0303-00503}{Mahasiswa mengerjakan UAS yang mencakup vocabulary dan grammatical functions, pemahaman bacaan/simakan, serta penulisan karangan pendek berdasarkan Topic 4--8 (Computer Network, Websites, Careers in IT, IT Support Staff, Workstation Health and Safety).}{Mengerjakan soal secara individu dalam waktu yang ditentukan: soal vocabulary-grammar, reading/listening comprehension, dan menulis karangan pendek sesuai instruksi.}{Lembar jawaban UAS (tulis) dan/atau rekaman respon lisan.}{Ketepatan vocabulary dan grammar & 5 & Lebih dari 80\% soal vocabulary dan grammar (if-clause, modals, prepositions, expressions of prohibition) dijawab benar. & 60--80\% soal dijawab benar; kesalahan terkonsentrasi pada bentuk yang kompleks. & Kurang dari 60\% soal dijawab benar. \\ \\
Pemahaman bacaan/simakan & 5 & Menjawab pertanyaan comprehension dengan benar, menemukan informasi rinci, dan menyimpulkan isi teks secara tepat. & Sebagian besar pertanyaan dijawab benar, tetapi kesimpulan atau informasi rinci kurang tepat. & Jawaban comprehension banyak yang keliru atau tidak menunjukkan pemahaman isi teks. \\ \\
Kualitas karangan & 5 & Karangan sesuai perintah, terorganisasi runtut, menggunakan grammatical functions target dengan benar, dan istilah TI yang tepat. & Karangan menjawab perintah, tetapi organisasi atau ketepatan grammar-vocabulary masih lemah di beberapa bagian. & Karangan tidak sesuai perintah, tidak runtut, atau kesalahan bahasa mendominasi. \\}

\begin{tblr}{width=\textwidth,colspec={|Q[l,m,3.2cm]|X[l,m]|},hlines,vlines,hspan=minimal,row{1-3}={bg=headgray},rowsep=1.5pt,colsep=3pt}
Jadwal Pelaksanaan: & Minggu 1--17 sesuai RPS. Setiap evaluasi memiliki RTM dan rubrik multi-kriteria. \\
Lain-Lain Yang Diperlukan: & LMS, perangkat presentasi, audio files, speakers, dan aplikasi pendukung (word processor, video editor, Grammarly/spelling checker). \\
Pustaka: & Mengacu pustaka utama dan pendukung pada bagian RPS. \\
\end{tblr}
